% Chapter 19 -- Claim: Most Published Anomalies Fail to Replicate
% Requires: amsmath, amssymb, booktabs, array, tabularx, longtable, graphicx, tcolorbox (breakable, skins), hyperref or url.
% Defines the callout box only if the preamble has not.
\providecolor{fenavy}{HTML}{1B2A6B}
\providecolor{feshade}{HTML}{F3F5FA}
\makeatletter
\@ifundefined{callout}{%
  \newtcolorbox{callout}{breakable, enhanced,
    colback=white, colframe=gray!70, boxrule=0.5pt, arc=1mm,
    borderline west={2.5pt}{0pt}{gray!70},
    left=8pt, right=8pt, top=4pt, bottom=4pt}%
}{}
\makeatother

\providecolor{fenavy}{HTML}{1B2A6B}
\providecolor{feshade}{HTML}{F3F5FA}
\makeatletter
% Shaded box for optional technical material that general readers may skip.
\@ifundefined{technote}{%
  \newtcolorbox{technote}{breakable, enhanced,
    colback=feshade, colframe=fenavy!60, boxrule=0.5pt, arc=1mm,
    left=8pt, right=8pt, top=4pt, bottom=4pt}%
}{}
\makeatother

\chapter{Claim: Most Published Anomalies Fail to Replicate}
\label{ch:19}

News stories regularly report a new study showing that some simple rule, such as buying cheap stocks or stocks with rising profits, has beaten the market. Other sciences, such as psychology and medicine, have discovered in recent years that many published findings could not be repeated when others tried. This chapter asks whether finance has the same problem, and how a reader can tell a lasting discovery from a lucky one.

By the 2010s academic finance had published hundreds of variables said to predict which stocks would earn more than others. Harvey, Liu and Zhu (2016) counted more than 300. They argued that, because so many had been tried, a new factor should clear a $t$-statistic of about 3 rather than 2. The worry was not new. Chapter~\ref{ch:01} showed how searching produces records that luck alone can explain. What was new was that the worry now applied to a whole literature. McLean and Pontiff (2016) took a different approach. Pure data mining predicts that returns vanish once the original sample ends, while mispricing predicts that they shrink after publication, as investors trade on them. They followed 97 predictors after their original samples ended and after they were published. Average returns were about 26\% lower after the sample ended, and 58\% lower after publication.

Two large replication studies sharpened the debate. Hou, Xue and Zhang (2020) asked how many anomalies survive when tiny stocks count for little, and rebuilt 452 anomalies. They weighted stocks by market value and set group cut-offs using NYSE stocks only, two choices that limit the influence of the tiniest stocks, called micro-caps. About two thirds of the anomalies then failed even the conventional test. Jensen, Kelly and Pedersen (2023) asked whether a verdict that uses the evidence from all factors at once would be kinder, and rebuilt 153 factors in 93 countries. They used a Bayesian method that pulls each factor's estimate toward what similar factors suggest, and concluded that most factors do replicate. Because much of the disagreement lay in each study's code, Chen and Zimmermann (2022) published open code and data for more than 200 predictors, so that anyone could check. Respected researchers using largely the same data reached opposite conclusions, because they asked slightly different questions and applied different standards.

At stake is the credibility of a whole field. If most anomalies are false, much of the factor-investing industry rests on noise. If most are real but fade after publication, investors are learning from research and trading the patterns away. The verdict depends less on the data than on the rules of evidence, so this chapter asks the reader to choose those rules and defend them. The claim on trial\index{claims-on-trial method} is: \textbf{most published cross-sectional return predictors do not replicate once micro-caps, multiple testing and publication are accounted for.}

The lab in Section~\ref{sec:19-lab} uses the free Open Source Asset Pricing data of Chen and Zimmermann. It applies the standards of Hou, Xue and Zhang (2020) and of Jensen, Kelly and Pedersen (2023) to the same predictors, and it measures the decay found by McLean and Pontiff (2016). The goal is to see why careful researchers reached opposite conclusions, and which choices drive the difference.

By the end of the chapter you will be able to:

\begin{itemize}
  \item explain market efficiency\index{efficient market hypothesis}, the joint-hypothesis problem\index{joint-hypothesis problem}, and the kinds of replication, from reproduction\index{reproduction} to pure, scientific and out-of-sample replication\index{replication!pure, scientific and out-of-sample};
  \item explain how publication bias\index{publication bias} inflates published results, and how empirical Bayes\index{empirical Bayes} shrinkage\index{shrinkage} corrects noisy estimates;
  \item set out the prosecution's and the defense's cases, and the design choices that lead them to opposite verdicts;
  \item separate post-publication\index{anomaly!post-publication decay} decay caused by selection from decay caused by investors' learning, and defend your own verdict.
\end{itemize}


\section{Foundations: market efficiency, anomalies and replication}
\label{sec:19-foundations}

\subsection{The efficient market hypothesis}

Section~\ref{sec:01-efficient-markets} introduced the efficient market hypothesis and Fama's (1970) three forms, weak, semi-strong and strong. The debate in this chapter concerns the semi-strong form\index{efficient market hypothesis!semi-strong form}: whether publicly available characteristics of firms, such as their valuation ratios or past returns, predict returns beyond what risk justifies.

\subsection{The joint-hypothesis problem}

Section~\ref{sec:01-efficient-markets} also introduced the joint-hypothesis problem. A strategy that beats a model such as the CAPM\index{capital asset pricing model (CAPM)} shows either that the market is inefficient or that the model is wrong, and the two cannot be separated (Fama 1970). For anomalies this has a direct consequence. One researcher can read the same evidence as mispricing\index{mispricing}, and another as a reward for a risk the model leaves out.

\subsection{Anomalies}

An \textbf{anomaly}\index{anomaly|textbf} is a pattern in returns that a benchmark asset pricing model cannot explain. It is a characteristic, such as a valuation ratio or a past return, that predicts future returns after adjusting for the model's risk factors. Anomalies are usually found with portfolio sorts\index{portfolio sort} and long--short returns (Section~\ref{sec:05-foundations}). They are judged by the t-statistic\index{t-statistic} of the average return or of the alpha\index{alpha}.

\subsection{Replication, reproduction and robustness}

Terms are used loosely, but a useful distinction is:

\begin{itemize}
  \item \textbf{Reproduction}\index{reproduction|textbf}: the same data and code give the same numbers. It checks for errors.
  \item \textbf{Pure replication}\index{replication!pure, scientific and out-of-sample|textbf}: the same method on the same kind of data, perhaps rebuilt independently, gives the same result. It checks that the finding is not an artifact of one implementation.
  \item \textbf{Scientific replication} or \textbf{robustness}\index{robustness|textbf}: different but reasonable choices of method, sample or market give a similar result. It checks that the finding is about the world, not about the choices.
  \item \textbf{Out-of-sample replication}: the result holds in data not used to find it, such as later years or other countries.
\end{itemize}

The two sides of this chapter's debate differ largely in which of these they demand.

\subsection{Publication bias}

Journals are more likely to publish significant results. If only strong in-sample\index{in-sample versus out-of-sample} results are published, the published estimates are too high on average, even if every study was done honestly. The reason is that the published set has been selected for luck as well as for truth. Part E of this chapter's lab measures how large this effect can be.

\subsection{Shrinkage}

Suppose many noisy estimates are made at once. Each one can be improved by pulling it toward the average of all of them. Noisy estimates are pulled more, and precise ones less. A very high estimate is more likely to be partly luck than a moderate one, so pulling it back makes it more accurate on average. This idea comes from \textbf{empirical Bayes}\index{empirical Bayes|textbf} statistics (Efron and Morris 1975). It is called ``empirical'' because the data themselves supply the average toward which the estimates are pulled. It is the basis of the defense's method.

\textbf{Worked example.} One hundred honest researchers each test a different trading signal, and none of the signals has any real predictive power. By chance, about five will find t-statistics beyond $\pm 1.96$, and those five publish (Figure~\ref{fig:19-pubbias}). Every published paper reports a significant result, yet none of the results is real. A reader who saw only the journals would find a literature with a 100\% success rate.

\begin{figure}[htbp]
\centering
\scalebox{0.8}{%
\begin{tikzpicture}[declare function={npdf(\x) = exp(-\x*\x/2)/sqrt(2*pi);}]
\begin{axis}[
  width=0.68\linewidth, height=0.34\linewidth,
  axis lines=left, axis line style={-stealth},
  xmin=-4, xmax=4.1, ymin=0, ymax=0.66,
  xtick={-1.96,0,1.96}, xticklabels={$-1.96$,$0$,$1.96$},
  ytick=\empty,
  xlabel={t-statistic of a signal with no true power},
  ylabel={Share of signals},
  samples=200, domain=-4:4, clip=false,
]
  \addplot[name path=curve, draw=none] {npdf(x)};
  \path[name path=base] (axis cs:-4,0) -- (axis cs:4,0);
  \addplot[fenavy] fill between[of=curve and base, soft clip={domain=-4:-1.96}];
  \addplot[fenavy] fill between[of=curve and base, soft clip={domain=1.96:4}];
  \addplot[fegray, thick] {npdf(x)};

  \node[align=center] at (axis cs:0,0.56) {about 95 signals: not significant,\\never published};
  \node[fenavy, align=center, font=\footnotesize, anchor=south] (L) at (axis cs:-2.75,0.13) {about 2.5 signals:\\``significant,''\\published};
  \node[fenavy, align=center, font=\footnotesize, anchor=south] (R) at (axis cs:2.75,0.13) {about 2.5 signals:\\``significant,''\\published};
  \draw[-stealth, fenavy] (axis cs:-2.7,0.12) -- (axis cs:-2.3,0.03);
  \draw[-stealth, fenavy] (axis cs:2.7,0.12) -- (axis cs:2.3,0.03);
\end{axis}
\end{tikzpicture}%
}
\caption{If only significant results are published, the journals show only the lucky draws. Publication bias. One hundred signals with no true power produce t-statistics spread around zero. Only the few in the shaded tails, about 2.5 on each side, look significant and get published. The published record shows only the tails.}
\label{fig:19-pubbias}
\end{figure}

\section{How it works: the factor zoo and how replications are designed}
\label{sec:19-how}

\textbf{The zoo.} For decades, finance explained why some stocks earn more than others with a single factor, the market, as in the CAPM. Later it used a handful, such as size, value and momentum\index{momentum}. By 2011 the literature had produced so many candidate predictors that John Cochrane, in his presidential address to the American Finance Association, called it a ``zoo of new factors'' (Cochrane 2011). He argued that the field needed to sort out which factors were truly independent, which were redundant, and which were simply noise. Harvey, Liu and Zhu (2016) counted 316 factors published in leading journals, and the count has kept growing.

The problem is statistical as much as economic. Thousands of researchers test thousands of variables on the same U.S. stock data, so some variables will look significant by luck alone. At the usual cut-off of a t-statistic near 2, about one useless variable in twenty passes, and journals mostly publish the ones that pass (Section~\ref{sec:19-foundations}). The published zoo is therefore a mixture of real effects and false discoveries, much as a stack of winning lottery tickets says little about whether anyone has a winning system. Many of its members are also the same idea measured in slightly different ways, such as several versions of the finding that cheap stocks beat expensive ones. Since Harvey, Liu and Zhu's count, two open data projects have made it possible to rebuild most of the zoo:

\begin{itemize}
  \item Chen and Zimmermann's \textbf{Open Source Asset Pricing}\index{Open Source Asset Pricing|textbf} (2022) replicates more than 200 cross-sectional predictors from the original papers' descriptions and publishes their portfolio returns and code.
  \item Jensen, Kelly and Pedersen (2023) built a \textbf{global factor database} of 153 factors across 93 countries, with harmonized definitions, meaning each factor is built the same way in every country.
\end{itemize}

No table can list hundreds of predictors, but none needs to, because most of them are variations on a few ideas. Hou, Xue and Zhang\index{q-factor model} (2020) sort their 452 anomalies into six categories. Table~\ref{tab:19-zoo} shows two or three well-known predictors from each. The sign shows whether stocks with a high value of the predictor earned higher (+) or lower ($-$) average returns in the original study. The Open Source Asset Pricing and JKP documentation give the full lists, with definitions and code for every predictor. Jensen, Kelly and Pedersen (2023) group their 153 factors into 13 similar themes.

\begin{table}[htbp]
\centering
\small
\caption{A map of the factor zoo: representative predictors in the six categories of Hou, Xue and Zhang (2020).}
\label{tab:19-zoo}
\renewcommand{\arraystretch}{1.15}
\begin{tabularx}{\linewidth}{@{}>{\raggedright\arraybackslash}p{0.13\linewidth}>{\raggedright\arraybackslash}p{0.17\linewidth}>{\raggedright\arraybackslash}X>{\raggedright\arraybackslash}p{0.21\linewidth}c@{}}
\toprule
Category & Predictor & What it measures & Original study & Sign \\
\midrule
Momentum & Price momentum & Return over the past 12 months, skipping the most recent month & Jegadeesh and Titman (1993) & $+$ \\
 & Earnings surprise & Change in quarterly earnings from a year earlier, scaled by its volatility & Bernard and Thomas (1989) & $+$ \\
 & Short-term reversal & Return over the past month & Jegadeesh (1990) & $-$ \\
\addlinespace
Value versus growth & Book-to-market & Book value of equity divided by market value & Rosenberg, Reid and Lanstein (1985); Fama and French (1992) & $+$ \\
 & Earnings-to-price & Earnings divided by market value & Basu (1977) & $+$ \\
\addlinespace
Investment & Asset growth & Annual percentage growth in total assets & Cooper, Gulen and Schill (2008) & $-$ \\
 & Net share issuance & Annual change in shares outstanding & Pontiff and Woodgate (2008) & $-$ \\
 & Accruals & Earnings minus operating cash flow, scaled by assets & Sloan (1996) & $-$ \\
\addlinespace
Profitability & Gross profitability & Revenue minus cost of goods sold, divided by assets & Novy-Marx (2013) & $+$ \\
 & F-score & Sum of nine accounting signals of financial strength & Piotroski (2000) & $+$ \\
\addlinespace
Intangibles & R\&D intensity & Research and development spending divided by market value & Chan, Lakonishok and Sougiannis (2001) & $+$ \\
\addlinespace
Trading frictions & Size & Market value of equity & Banz (1981) & $-$ \\
 & Idiosyncratic volatility & Volatility of returns not explained by the Fama--French factors & Ang, Hodrick, Xing and Zhang (2006) & $-$ \\
 & Illiquidity & Average absolute daily return per dollar of trading volume & Amihud (2002) & $+$ \\
 & Maximum daily return & Largest daily return in the past month & Bali, Cakici and Whitelaw (2011) & $-$ \\
\bottomrule
\end{tabularx}
\end{table}

\textbf{Design choices in a replication\index{replication}.} A replication of an anomaly must make many choices, each defensible, that can change the result:

\begin{center}
\small
\begin{tabularx}{0.9\linewidth}{@{}>{\raggedright\arraybackslash}p{0.17\linewidth}>{\raggedright\arraybackslash}X>{\raggedright\arraybackslash}p{0.34\linewidth}@{}}
\toprule
Choice & Options & Why it matters \\
\midrule
Universe & All stocks, or excluding micro-caps (for example below the NYSE 20th size percentile) & Many anomalies are concentrated in tiny stocks that are costly to trade and hold little value \\
Weighting & Equal-weighted or value-weighted portfolios & Equal weighting gives micro-caps large influence \\
Breakpoints & Deciles, quintiles or terciles; NYSE or all-stock breakpoints & Extreme portfolios magnify signals and noise \\
Sample & The original paper's sample, or the longest available & Later years are out of sample \\
Significance hurdle & $t > 1.96$, $t > 3$, a false-discovery-rate rule, or a Bayesian posterior & Determines what counts as a success \\
Benchmark & Raw returns, CAPM alpha, or alpha against a multifactor model & A factor explained by others may be real but redundant \\
\bottomrule
\end{tabularx}
\end{center}

\textbf{NYSE breakpoints}\index{NYSE breakpoints|textbf} set the cut-offs between groups using only stocks listed on the New York Stock Exchange. These stocks are larger on average, so the many tiny Nasdaq stocks do not crowd the extreme portfolios. A \textbf{Bayesian posterior}\index{Bayesian posterior|textbf} is the probability of a hypothesis after a prior belief has been combined with the data.

\textbf{Pure and scientific replication again.} A replication that copies the original method tests whether the paper's numbers can be reproduced. One that changes the universe, weighting or hurdle tests whether the finding is robust. Both are valuable, but they answer different questions, and the debate in this chapter turns partly on that difference.

\section{The claim}
\label{sec:19-claim}

Stated so that it can fail:

\begin{callout}
\textbf{Most published cross-sectional return predictors do not reflect a real, persistent relationship: they fail to replicate under reasonable standards, or their returns disappear outside the sample in which they were found.}
\end{callout}

``Reasonable standards'' is where the debate lies. The claim can be tested in three ways:

\begin{itemize}
  \item \textbf{Replication within the original sample.} Do the predictors' long--short returns clear a significance hurdle when rebuilt independently? Which hurdle, and with which portfolio construction?
  \item \textbf{Out of sample.}\index{in-sample versus out-of-sample|textbf} Do returns persist after the original sample ends and after publication?
  \item \textbf{Statistically corrected.} After accounting for the number of predictors tested and for publication bias, what share have a high probability of a positive true return?
\end{itemize}

\section{The prosecution}
\label{sec:19-prosecution}

\textbf{Harvey, Liu and Zhu (2016)} argued from statistics. Hundreds of factors have been tested, and many more were tested but never published. The usual hurdle of $t > 2$ therefore lets through too many false discoveries. After adjusting for the number of tests, they proposed that a new factor should clear a t-statistic of about 3. They concluded that many published factors were likely false.

\textbf{Hou, Xue and Zhang (2020)}\index{q-factor model|textbf} argued from replication. They rebuilt 452 anomalies with one common method, designed to limit the influence of micro-caps\index{micro-caps}: NYSE breakpoints and value-weighted\index{index weighting!value-weighted} returns. With that method, 65\% of the anomalies failed the single-test hurdle of an absolute t-statistic of 1.96. Raising the hurdle to 2.78, to allow for multiple testing, raised the failure rate to 82\%. Anomalies based on \textbf{trading frictions}\index{trading frictions|textbf}, the costs and obstacles of trading, failed most often: 96\% of them failed. This group includes illiquidity\index{illiquidity} and measures of idiosyncratic risk. \textbf{Idiosyncratic risk}\index{risk!idiosyncratic|textbf} is the part of a stock's volatility\index{volatility} that the market and other factors do not explain. Their conclusion was that most anomalies are either false or confined to micro-caps, which are too small and costly to matter.

The prosecution's case rests on two premises: that the right hurdle is high because many tests have been run, and that the right portfolio construction gives little weight to the smallest stocks.

\section{The defense}
\label{sec:19-defense}

\textbf{Jensen, Kelly and Pedersen (2023)} argued that the replication crisis was overstated. Using a Bayesian model of factor replication, they found that:

\begin{itemize}
  \item most factors can be replicated;
  \item they cluster into 13 themes, most of which are significant parts of the optimal portfolio;
  \item they work out of sample in a large new data set covering 93 countries;
  \item the evidence is strengthened rather than weakened by the large number of factors, because related factors corroborate one another.
\end{itemize}

Their approach treats each factor's alpha as a draw from a common distribution and shrinks noisy estimates toward the average of related factors, rather than judging each one in isolation against a fixed hurdle.

\textbf{Chen and Zimmermann (2022)} asked whether published results can be reproduced at all, and found that almost all predictors whose original papers reported clear statistical evidence can be reproduced with the original methods. \textbf{Chen and Zimmermann (2020)} asked how much journals' preference for strong results inflates published returns, and estimated publication bias directly across 156 published predictors. They concluded that returns adjusted for bias are only about 12\% smaller than in-sample returns. For predictors that survive journal review, they found that a t-statistic hurdle of about 1.8 is enough to control for multiple testing. \textbf{Chen (2021)} asked whether \textbf{p-hacking}\index{p-hacking|textbf}, that is, trying analyses until one crosses the significance line, could explain the published evidence, and argued that published t-statistics are too large to be explained by p-hacking alone.

The defense's case rests on two different premises. First, the question is whether the original findings are real, not whether they survive a method that plays down small stocks. Second, evidence should be pooled across related factors rather than judged one test at a time.

\section{Post-publication decay}
\label{sec:19-decay}

Whatever the replication standard, performance outside the original sample is a separate question. \textbf{McLean and Pontiff (2016)}\index{anomaly!post-publication decay|textbf} followed 97 predictors after their original samples ended and after publication. Portfolio returns were 26\% lower after the sample ended and 58\% lower after publication (Figure~\ref{fig:19-decay}). They read the first figure as an upper bound on the effect of data mining\index{data mining}. Nothing in the world changed when a sample ended, except that researchers stopped looking. They put the further decline after publication down to investors trading on the published results.

\begin{figure}[htbp]
\centering
\scalebox{0.8}{%
\begin{tikzpicture}
\begin{axis}[
  width=0.8\linewidth, height=0.34\linewidth,
  ybar, bar width=0.68cm,
  axis lines=left, axis line style={-stealth},
  ymin=0, ymax=118, ytick={0,50,100},
  ylabel={Return (in sample $= 100$)},
  symbolic x coords={A,B,C},
  xtick=data,
  xticklabels={In sample, {After sample,\\before\\publication}, {After\\publication}},
  xticklabel style={align=center, text width=2.3cm, font=\small},
  enlarge x limits=0.25,
  nodes near coords,
]
  \addplot[fill=fenavy, draw=none] coordinates {(A,100) (B,74) (C,42)};
\end{axis}
\end{tikzpicture}%
}
\caption{Published patterns lost more than half of their returns after publication. Post-publication decay in McLean and Pontiff (2016). With in-sample returns set to 100, the average predictor earned 74 between the end of its original sample and publication, and 42 after publication.}
\label{fig:19-decay}
\end{figure}

The two readings matter for the claim. If most of the decline happens when the sample ends, the original results were inflated by selection, which is a sign of data mining. If most happens after publication, the results were real but investors have traded them away. That is a sign of learning, and it suggests the original finding was true. McLean and Pontiff's numbers suggest that both forces matter, and that the publication effect is the larger one.

This chapter's lab adds a check the original study could not make: a simulated literature in which the truth is known. If only significant results are published, returns appear to fall after the sample ends even when nothing changes. The size of that purely statistical decline is a benchmark for how much of the observed decay is selection.

\section{Why the two sides disagree}
\label{sec:19-why}

The prosecution and the defense agree on most facts. They disagree on four choices.

\begin{center}
\small
\begin{tabularx}{0.9\linewidth}{@{}>{\raggedright\arraybackslash}p{0.17\linewidth}>{\raggedright\arraybackslash}X>{\raggedright\arraybackslash}X@{}}
\toprule
Choice & Prosecution (Hou, Xue and Zhang; Harvey, Liu and Zhu) & Defense (Jensen, Kelly and Pedersen; Chen and Zimmermann) \\
\midrule
Which stocks count & De-emphasize micro-caps: NYSE breakpoints, value weights & Replicate the original construction; micro-caps are part of the market \\
Hurdle & High: $t$ of about 3 or 2.78, for multiple testing & Lower, or a Bayesian posterior; multiple-testing corrections for published factors need not be severe \\
Unit of evidence & Each factor judged on its own & Related factors pooled; corroboration strengthens the evidence \\
Kind of replication & Scientific: does it survive reasonable alternative choices? & Pure: are the original findings reproducible and real? \\
\bottomrule
\end{tabularx}
\end{center}

Neither side is wrong on its own terms. A factor that appears only among micro-caps may be real, yet irrelevant to an investor who cannot trade micro-caps. A factor that fails a $t > 3$ hurdle may be real but measured with a lot of noise. A Bayesian estimate that pools related factors is more accurate if its assumptions are right. It is misleading if a whole theme was mined together.

The practical lesson is that ``does this anomaly replicate?'' is not a well-posed question until the purpose is stated. For understanding how markets price risk, pure replication and pooled evidence are appropriate. For deciding what to invest in, value-weighted, liquid-universe, out-of-sample, after-cost evidence is what matters.

\section{Lab~\thechapter: Replicating both sides of the anomaly debate}
\label{sec:19-lab}

\begin{callout}
\textbf{Lab~\ref{ch:19} at a glance.} \emph{Papers replicated:} Hou, Xue and Zhang (2020); Jensen, Kelly and Pedersen (2023); McLean and Pontiff (2016). All data are free. \emph{The question:} applying each side's standards to the same data, what share of published predictors survive, and how much of their return remains after their samples end and after publication? \emph{Data:} Open Source Asset Pricing: monthly long--short returns for more than 200 predictors as constructed in the original papers, and the signal documentation giving each paper's sample period and publication year. Optionally, a second construction of the same signals, value-weighted or excluding micro-caps, for the methodology comparison. \emph{Code:} repository folder \texttt{ch10-anomalies/}, with a guided notebook.
\end{callout}

\textbf{The setup.}

\begin{itemize}
  \item \textbf{Part A, replication.} Each predictor's t-statistic over its own original sample. The share clearing $t > 1.96$, $t > 3$\index{t-statistic!$t>3$ hurdle}, and a Benjamini--Hochberg test, which limits the expected share of false discoveries among the predictors it declares significant.
  \item \textbf{Part B, decay.} Average returns in-sample, after the sample ended but before publication, and after publication. A \textbf{pooled regression}\index{regression!pooled|textbf}, which combines all predictors into one regression in the style of McLean and Pontiff, measures the proportional decline in each period with errors clustered by\index{standard error!clustered} month.
  \item \textbf{Part C, methodology.} The same statistics for an alternative construction, if supplied. This is the core of the disagreement between Hou, Xue and Zhang and the defense.
  \item \textbf{Part D, shrinkage.} An empirical Bayes estimate of each predictor's true mean, pulled toward the average in proportion to its noise; the share with more than 95\% probability of a positive true mean; and whether shrunk estimates predict post-publication returns better than raw ones.
  \item \textbf{Part E, known truth.} A simulated literature in which only signals with in-sample $t > 2$ are published, run with no true decay and with a 30\% decay after publication.
\end{itemize}

\textbf{What readers will find.} The simulation alone delivers one of the chapter's central lessons. The lab first ran a simulated literature with \textbf{no} true change after publication. The pooled regression still measured returns about 45\% lower after the samples ended and about 54\% lower after publication. This decline came purely from publication selecting lucky in-sample results. When a true 30\% post-publication decay was added, the measured post-publication decline rose to about 68\%. Readers who run the lab on real data can compare McLean and Pontiff's figures, and their own, with this benchmark.

\subsection*{Your verdict}

Reach your own verdict on the claim from the numbers in your own run, using the scale of Section~\ref{sec:01-method}. Then compare it with the published evidence summarized in Section~\ref{sec:19-verdict}. If your verdict differs, explain which of your results led you there. State the standard of evidence you apply and why, and date your verdict with the sample you tested.

\section{Where the evidence stands}
\label{sec:19-verdict}

The published studies cited in this chapter point to a split picture: \textbf{most published anomalies are not pure statistical flukes. But most are smaller than reported, concentrated in small stocks, partly traded away, and not profitable to trade after costs.} The table summarizes where the evidence supports the claim and where it does not, as of November 2026, and names the source of each finding. The lab lets you verify these findings with free data.

{\small
\begin{longtable}{@{}>{\raggedright\arraybackslash}p{0.19\linewidth}>{\raggedright\arraybackslash}p{0.10\linewidth}>{\raggedright\arraybackslash}p{\dimexpr 0.48\linewidth-6\tabcolsep\relax}>{\raggedright\arraybackslash}p{0.23\linewidth}@{}}
\toprule
Situation & Holds? & Why & Evidence \\
\midrule
\endfirsthead
\multicolumn{4}{@{}l}{\emph{(continued)}} \\
\toprule
Situation & Holds? & Why & Evidence \\
\midrule
\endhead
\midrule
\multicolumn{4}{r@{}}{\emph{(continued on the next page)}} \\
\endfoot
\bottomrule
\endlastfoot
Original construction, conventional hurdle & \textbf{No} & Predictors whose original papers reported clear evidence almost all reproduce with the original methods. & Chen and Zimmermann (2022) \\
\addlinespace[3pt]
Micro-caps de-emphasized, value-weighted & \textbf{Yes} & 65\% of 452 anomalies fail to clear $t = 1.96$, and 82\% fail a hurdle of 2.78; 96\% of the trading-frictions anomalies fail. & Hou, Xue and Zhang (2020) \\
\addlinespace[3pt]
The right hurdle for multiple testing & \textbf{Disputed} & One side asks for a $t$ of about 3; the other estimates that about 1.8 is enough for published predictors. & Harvey, Liu and Zhu (2016); Chen and Zimmermann (2020) \\
\addlinespace[3pt]
Bayesian pooling across factors and countries & \textbf{No} & Most factors replicate, cluster into 13 themes and work out of sample in 93 countries. & Jensen, Kelly and Pedersen (2023) \\
\addlinespace[3pt]
Out of sample, after publication & \textbf{Partly} & Returns were 26\% lower after the original samples and 58\% lower after publication, and part of the decline is selection. & McLean and Pontiff (2016); Part E of the lab \\
\addlinespace[3pt]
After costs, for an investor & \textbf{Yes} & Little expected return remains for the average anomaly. & Chen and Velikov (2023); Chapter~\ref{ch:17} \\
\end{longtable}
}

Whether these findings make most anomalies ``false'' depends on whether the question is scientific or practical. The two sides agree on most facts and disagree on the four choices in Section~\ref{sec:19-why}, so a verdict on this claim is a choice of standard as much as a reading of the data.

\textbf{What would change the picture.} A shared replication standard, such as pre-registered\index{pre-registration} hurdles (standards fixed before the data are examined) or value-weighted, micro-cap-robust results reported alongside the original construction, would narrow the disagreement more than any new data set.

\subsection{Open questions}
\label{sec:19-open}

\begin{itemize}
  \item \textbf{What replication standard should the field adopt?} Should journals require value-weighted, micro-cap-robust results, pre-registered\index{pre-registration} hurdles, or Bayesian posteriors?
  \item \textbf{How many signals were tested but never published?} The unknown size of this \textbf{file drawer}\index{publication bias|textbf}, the studies left unpublished, limits every correction for multiple testing.
  \item \textbf{Is post-publication decay learning or selection?} Part E's benchmark separates the two in simulation; doing so convincingly in real data remains open.
  \item \textbf{Do anomalies reflect risk or mispricing?} Their decay after publication favors mispricing for many; their persistence\index{persistence} across countries favors something more durable for a few themes.
\end{itemize}

\section{What practitioners and investors can use}
\label{sec:19-investors}

The evidence in Section~\ref{sec:19-verdict} suggests several practical lessons.

\textbf{A checklist for judging any new factor fund or signal:}

\begin{enumerate}
  \item \textbf{Is the original evidence strong?} Ask for the t-statistic in the original study and whether it clears about 3.
  \item \textbf{Does it hold among large, liquid stocks, value-weighted?} If the effect lives in micro-caps, it is unlikely to be investable at scale.
  \item \textbf{What happened after publication?} Ask for returns from the paper's publication date forward. Expect roughly half the backtest\index{backtest} at best.
  \item \textbf{Does it belong to a theme that works?} A signal related to established themes such as value, momentum, profitability or low risk has corroborating evidence; an isolated signal has none.
  \item \textbf{What survives costs?} Ask for turnover\index{turnover} and net returns at realistic costs (Chapter~\ref{ch:17}).
  \item \textbf{How many signals did the provider test before choosing this one?} Apply this chapter's multiple-testing corrections to the answer.
\end{enumerate}

\begin{callout}
\textbf{What Failed: the trading-frictions anomalies.} A family of anomalies based on liquidity\index{liquidity}, idiosyncratic volatility\index{idiosyncratic volatility} and other trading frictions looked strong in the papers that introduced them. Hou, Xue and Zhang played down micro-caps with NYSE breakpoints and value weights. In their replication, 96\% of the anomalies in this group then failed to clear a t-statistic of 1.96. Their returns had come mostly from the smallest and least liquid stocks, where costs and frictions are largest. The claim that frictions-based anomalies offer returns an investor can capture \textbf{fails}. Whether they reveal something real about how illiquid stocks are priced is a separate, open question.
\end{callout}
